% ============================================================================
%  Chapter 4 -- Representing Digital Information                 (Lecture 04)
% ============================================================================
\chapter{Representing Digital Information}
\label{ch:information}

\epigraphline{\texttt{printf("Hello, world!");}}
  {Brian Kernighan and Dennis Ritchie, \emph{The C Programming Language}, 1978}

\begin{chapgoals}
By the end of this chapter you will be able to
\begin{itemize}
\item explain the difference between analog and digital information;
\item read a pattern of bits as a natural number, in binary and in hexadecimal;
\item decode a \prg colour such as \hex{07E0} into its red, green and blue fields;
\item say how characters and strings are stored in memory, in ASCII and in UTF-8;
\item describe how an array lives in memory and how assembly reaches one element;
\item tell a deterministic algorithm from a probabilistic one, and place machine learning among them.
\end{itemize}
\end{chapgoals}

In Chapter~\ref{ch:firstcontact} the reader typed \hex{07e0} and obtained
green, typed 255 and obtained a bar that starts again, typed \code{.asciz} and
obtained a greeting, typed \code{.word} and obtained a counter. This chapter
gives those things their names. Its thesis fits in one sentence, which the
rest of the chapter illustrates four times over: \emph{everything in the
machine is a pattern of bits, and the meaning of a pattern is a convention}.
The two chapters that follow treat numbers formally.

% ----------------------------------------------------------------------------
\section{Information}

\subsection{The flow of data}

Every use of a computer has the same shape. \emph{Data} enter from the world:
a key, a sensor, a file. A \emph{program}, the embodiment of an algorithm,
transforms them. \emph{Results} leave: a picture, a sound, a number. And
somebody receives \emph{knowledge}, that is, information on which it is
possible to act. On \prg the four stages are, in order, the function
\code{prg32\_input\_read}, the programmer's \code{update}, the programmer's
\code{draw}, and the player who reacts to the screen.

The last stage has, until now, required a person. Artificial intelligence, to
which the end of this chapter returns, is an intelligence of its own kind:
powerful pattern recognition, without self-awareness and without subjective
experience.

\subsection{Analog and digital information}

\begin{definition}{Analog and digital}
An \term{analog} quantity varies continuously: a voltage, a sound pressure,
the hands of a watch. A \term{digital} quantity is a sequence of symbols
taken from a finite alphabet: the digits of a display.
\end{definition}

To \term{digitise} an analog quantity means to do two things to it
(Figure~\ref{fig:info:sampling}): to \term{sample} it in time, reading its
value at regular instants, and to \term{quantise} it in value, replacing each
reading by the nearest of a finite number of levels.

\begin{figure}[htb]
\centering
\begin{tikzpicture}[x=0.75cm,y=0.42cm]
  \draw[->,prgink] (0,0) -- (12.6,0) node[right,lbl]{time};
  \draw[->,prgink] (0,0) -- (0,8.6) node[above,lbl]{value};
  \foreach \l in {1,...,7} \draw[prgrule,thin] (0,\l) -- (12.2,\l) node[right,font=\ttfamily\scriptsize,text=prgcomment]{\l};
  \draw[prgaccent,very thick,domain=0:12,smooth,samples=80] plot (\x,{4+2.6*sin(\x*38)+0.9*sin(\x*95)});
  \foreach \x/\y in {0/4,1/6,2/6,3/6,4/6,5/4,6/2,7/1,8/2,9/4,10/4,11/6,12/7}{
    \draw[prgamber,thick] (\x,0) -- (\x,\y); \fill[prgamber] (\x,\y) circle (2.2pt);
  }
  \node[lbl,text=prgaccent] at (3.8,8.4) {analog signal};
  \node[lbl,text=prgamber] at (8.2,-0.9) {samples, quantised to 3 bits: 4, 6, 6, 6, 6, 4, 2, 1, \ldots};
\end{tikzpicture}
\caption{Digitising: sampling in time and quantising in value.}
\label{fig:info:sampling}
\end{figure}

Something is lost in both steps, and something decisive is gained. A finite
number of levels can be \emph{counted}, \emph{copied exactly} and
\emph{checked for errors}; a continuous quantity can only be measured, and
every copy degrades it. With two levels the symbol is the \term{bit}, which
is 0 or 1. Two levels are easy to build reliably with transistors, which are
off or on (Chapter~\ref{ch:electronics}), and that is why computers are
binary. Boolean algebra (Chapter~\ref{ch:boole1}) is the mathematics of two
levels.

\subsection{Bit, byte, word}

Bits are handled in groups of standard sizes (Table~\ref{tab:info:units}).
The fundamental fact about them is a count.

\begin{theorem}{Number of patterns}
With $n$ bits it is possible to write exactly $2^{n}$ distinct patterns.
\end{theorem}

\noindent\emph{Proof.} By induction on $n$. One bit has $2=2^{1}$ patterns.
If $n$ bits have $2^{n}$ patterns, each of them can be followed by a 0 or by a
1, which gives $2\cdot 2^{n}=2^{n+1}$ patterns of $n+1$ bits. \hfill$\square$

\begin{table}[htb]
\centering\small
\caption{Units of information.}\label{tab:info:units}
\begin{tabularx}{\textwidth}{@{}llrL@{}}
\toprule
\textbf{Unit} & \textbf{Size} & \textbf{Patterns} & \textbf{Where it was met}\\
\midrule
bit & 1 binary digit & 2 & one button of the joystick\\
nibble & 4 bits & 16 & one hexadecimal digit\\
byte & 8 bits & 256 & one ASCII character; the width of the bar (\code{andi 255})\\
halfword & 16 bits & 65\,536 & one RGB565 colour\\
word & 32 bits & 4\,294\,967\,296 & one \riscv register; one instruction; \code{.word}\\
KiB, MiB, GiB & $2^{10}$, $2^{20}$, $2^{30}$ bytes & -- & memory sizes: 1 KiB = 1024 bytes\\
\bottomrule
\end{tabularx}
\end{table}

\begin{prgconcept}
The meaning of a pattern is a convention, not a property of the bits. What a
pattern means is not written in it: it is an agreement between whoever writes
it and whoever reads it.
\end{prgconcept}

The size of a word depends on the architecture: on RV32 a word is 32 bits by
definition.

\subsection{Counting in base 2 and in base 16}

Table~\ref{tab:info:hex} lists the sixteen patterns of four bits with their
values in decimal and their hexadecimal digits. The four columns of a binary
number weigh 8, 4, 2 and~1. This table should be learnt by heart: it is the
multiplication table of this book.

\begin{table}[htb]
\centering\small
\caption{The sixteen patterns of four bits.}\label{tab:info:hex}
\begin{tabular}{@{}ccc@{\qquad\qquad}ccc@{}}
\toprule
\textbf{Dec} & \textbf{Binary} & \textbf{Hex} & \textbf{Dec} & \textbf{Binary} & \textbf{Hex}\\
\midrule
0 & \texttt{0000} & \texttt{0} & 8 & \texttt{1000} & \texttt{8}\\
1 & \texttt{0001} & \texttt{1} & 9 & \texttt{1001} & \texttt{9}\\
2 & \texttt{0010} & \texttt{2} & 10 & \texttt{1010} & \texttt{A}\\
3 & \texttt{0011} & \texttt{3} & 11 & \texttt{1011} & \texttt{B}\\
4 & \texttt{0100} & \texttt{4} & 12 & \texttt{1100} & \texttt{C}\\
5 & \texttt{0101} & \texttt{5} & 13 & \texttt{1101} & \texttt{D}\\
6 & \texttt{0110} & \texttt{6} & 14 & \texttt{1110} & \texttt{E}\\
7 & \texttt{0111} & \texttt{7} & 15 & \texttt{1111} & \texttt{F}\\
\bottomrule
\end{tabular}
\end{table}

\subsection{Positional notation: what \texttt{11011001} means}

In a \term{positional} system the value of a digit depends on its position:
each position weighs a power of the base. Base ten uses powers of ten, base
two powers of two, base sixteen powers of sixteen. The rule is the same; only
the base changes. The byte \texttt{11011001} is read as follows.

\begin{center}\small
\begin{tabular}{r|*{8}{>{\centering\arraybackslash}p{9mm}|}}
\multicolumn{1}{r}{weight} & \multicolumn{1}{c}{$2^7$} & \multicolumn{1}{c}{$2^6$} & \multicolumn{1}{c}{$2^5$} & \multicolumn{1}{c}{$2^4$} & \multicolumn{1}{c}{$2^3$} & \multicolumn{1}{c}{$2^2$} & \multicolumn{1}{c}{$2^1$} & \multicolumn{1}{c}{$2^0$}\\
\multicolumn{1}{r}{} & \multicolumn{1}{c}{\scriptsize 128} & \multicolumn{1}{c}{\scriptsize 64} & \multicolumn{1}{c}{\scriptsize 32} & \multicolumn{1}{c}{\scriptsize 16} & \multicolumn{1}{c}{\scriptsize 8} & \multicolumn{1}{c}{\scriptsize 4} & \multicolumn{1}{c}{\scriptsize 2} & \multicolumn{1}{c}{\scriptsize 1}\\\cline{2-9}
bit & \textbf{1} & \textbf{1} & 0 & \textbf{1} & \textbf{1} & 0 & 0 & \textbf{1}\\\cline{2-9}
\multicolumn{1}{r}{contributes} & \multicolumn{1}{c}{128} & \multicolumn{1}{c}{64} & \multicolumn{1}{c}{0} & \multicolumn{1}{c}{16} & \multicolumn{1}{c}{8} & \multicolumn{1}{c}{0} & \multicolumn{1}{c}{0} & \multicolumn{1}{c}{1}\\
\end{tabular}
\[128 + 64 + 16 + 8 + 1 = 217\]
\end{center}

Chapter~\ref{ch:integers} gives the general formula and the algorithms of
conversion. Here the aim is only to read a byte.

\subsection{Why programmers write hexadecimal}

One hexadecimal digit is exactly four bits, because $16 = 2^{4}$. Converting
between binary and hexadecimal is therefore a table lookup, digit by digit,
with no arithmetic:
\[
\underbrace{\texttt{0000}}_{\texttt{0}}\;\underbrace{\texttt{0111}}_{\texttt{7}}\;
\underbrace{\texttt{1110}}_{\texttt{E}}\;\underbrace{\texttt{0000}}_{\texttt{0}}
\qquad\Longrightarrow\qquad \hex{07E0}.
\]
A 32-bit register is eight hexadecimal digits: \hex{00130313} is readable,
thirty-two zeros and ones are not. The prefix \code{0x} marks hexadecimal in
assembly and in C, and the debugger of \prgqt shows the registers in this way.

The decimal value of \hex{07E0} is $7\cdot 256+14\cdot 16 = 2016$, and it is
of no use whatever: what matters in this number is the \emph{bit pattern}, as
the next section shows.

\subsection{A colour is a number: RGB565}

On \prg a colour is sixteen bits divided into three \term{fields}: five bits
of red, six of green, five of blue (Figure~\ref{fig:info:rgb}). The format is
called RGB565. The constant \code{.equ GREEN, 0x07e0} is the pattern in which
green is at its maximum, 63 of 63, and red and blue are at zero. The eye is
most sensitive to green, which is why green receives the extra bit.

\begin{figure}[htb]
\centering
\begin{bytefield}[bitwidth=0.05\textwidth,bitheight=8mm]{16}
\bitheader[endianness=big]{0,4,5,10,11,15}\\
\bitbox{5}[bgcolor=prgred!35]{\texttt{00000}} & \bitbox{6}[bgcolor=prggreen!35]{\texttt{111111}} & \bitbox{5}[bgcolor=blue!25]{\texttt{00000}}\\
\bitbox[]{5}{\footnotesize red: 5 bits\\0 of 31} & \bitbox[]{6}{\footnotesize green: 6 bits\\63 of 63} & \bitbox[]{5}{\footnotesize blue: 5 bits\\0 of 31}
\end{bytefield}

\medskip
\begin{tikzpicture}[x=1.55cm]
  \foreach \i/\n/\v/\c in {0/BLACK/0000/{0,0,0}, 1/RED/F800/{248,0,0}, 2/GREEN/07E0/{0,252,0}, 3/BLUE/001F/{0,0,248}, 4/YELLOW/FFE0/{248,252,0}, 5/CYAN/07FF/{0,252,248}, 6/MAGENTA/F81F/{248,0,248}, 7/WHITE/FFFF/{255,255,255}}{
    \definecolor{sw}{RGB}{\c}
    \draw[prgsteel,fill=sw] (\i,0) rectangle ++(0.85,0.6);
    \node[font=\sffamily\scriptsize\bfseries] at (\i+0.42,-0.25) {\n};
    \node[font=\ttfamily\scriptsize,text=prgcomment] at (\i+0.42,-0.6) {0x\v};
  }
\end{tikzpicture}
\caption{RGB565: \hex{07E0} is green. Below, the colour constants of
\fname{prg32.h}.}
\label{fig:info:rgb}
\end{figure}

Sixteen bits, three fields. The same sixteen bits could be a number, two
characters or half an instruction: the meaning is in how the program uses
them.

\subsection{Numbers: naturals, integers, reals}

Numbers themselves are represented by three different conventions, which the
next two chapters develop.

\begin{itemize}
\item \textbf{Natural numbers} are the positional pattern itself: $n$ bits
  cover the interval from 0 to $2^{n}-1$.
\item \textbf{Signed integers} are written in \term{two's complement}: a
  32-bit register covers about $-2$ thousand million to $+2$ thousand million
  (Chapter~\ref{ch:integers}).
\item \textbf{Real numbers} are written in \term{fixed point} or in the
  \term{floating point} of the standard IEEE~754, where the value is
  $(-1)^{s}\times 1.f\times 2^{\,e-\mathit{bias}}$
  (Chapter~\ref{ch:reals}). Single precision uses 32 bits, gives about 7
  decimal digits and reaches about $3.4\times 10^{38}$; double precision uses
  64 bits, gives about 16 digits and reaches about $1.8\times 10^{308}$.
\end{itemize}

% ----------------------------------------------------------------------------
\section{Text}

Characters are numbers too.

\subsection{ASCII: one character, one byte}

The ASCII code (1963) assigns a number between 0 and 127, seven bits, to each
of 128 characters: the letters of English, the digits, punctuation and some
control codes (Table~\ref{tab:info:ascii}).

\begin{table}[htb]
\centering\small
\caption{Some ASCII characters.}\label{tab:info:ascii}
\begin{tabularx}{\textwidth}{@{}lrllL@{}}
\toprule
\textbf{Character} & \textbf{Decimal} & \textbf{Hex} & \textbf{Binary} & \textbf{Note}\\
\midrule
\texttt{A} & 65 & \hex{41} & \texttt{0100 0001} & capital letters start at 65\\
\texttt{a} & 97 & \hex{61} & \texttt{0110 0001} & the same as \texttt{A} with one more bit: $97-65=32$\\
\texttt{0} & 48 & \hex{30} & \texttt{0011 0000} & the character zero is not the number zero\\
space & 32 & \hex{20} & \texttt{0010 0000} & a character like any other\\
newline & 10 & \hex{0A} & \texttt{0000 1010} & a control character\\
NUL & 0 & \hex{00} & \texttt{0000 0000} & marks the end of a C string\\
\bottomrule
\end{tabularx}
\end{table}

The table was designed with care. Upper and lower case differ only in bit~5,
so that \code{xori} with \hex{20} switches the case of a letter: a first taste
of Boolean algebra applied to text. And the digit characters are consecutive,
so that subtracting 48 from the character of a digit gives the value of the
digit.

\subsection{Strings: a sequence of characters in memory}

A \term{string} is a sequence of characters stored in consecutive bytes. The
only question is how the reader of the string knows where it ends, and there
are two classical answers (Figure~\ref{fig:info:strings}). The \emph{C
string} marks the end with a byte of value zero. The \emph{Pascal string}
puts the length first.

\begin{figure}[htb]
\centering\small\ttfamily
\begin{tabular}{@{}r@{\quad}*{5}{|>{\centering\arraybackslash}p{7mm}}|@{\quad}l@{}}
\cline{2-6}
\textsf{C string} & C & I & A & O & \cellcolor{prgred!25}\textbackslash0 & \textsf{\footnotesize the end is marked by a zero byte}\\\cline{2-6}
\multicolumn{1}{r}{} & \multicolumn{1}{c}{\scriptsize 43} & \multicolumn{1}{c}{\scriptsize 49} & \multicolumn{1}{c}{\scriptsize 41} & \multicolumn{1}{c}{\scriptsize 4F} & \multicolumn{1}{c}{\scriptsize 00} & \\[6pt]
\cline{2-6}
\textsf{Pascal string} & \cellcolor{prgamber!25}4 & C & I & A & O & \textsf{\footnotesize the length comes first}\\\cline{2-6}
\multicolumn{1}{r}{} & \multicolumn{1}{c}{\scriptsize 04} & \multicolumn{1}{c}{\scriptsize 43} & \multicolumn{1}{c}{\scriptsize 49} & \multicolumn{1}{c}{\scriptsize 41} & \multicolumn{1}{c}{\scriptsize 4F} & \\
\end{tabular}

\bigskip
\begin{tabular}{@{}*{12}{|>{\centering\arraybackslash}p{5.2mm}}|@{}}
\hline
H & E & L & L & O & \textvisiblespace & P & R & G & 3 & 2 & \cellcolor{prgred!25}\textbackslash0\\\hline
\multicolumn{1}{c}{\scriptsize 48} & \multicolumn{1}{c}{\scriptsize 45} & \multicolumn{1}{c}{\scriptsize 4C} & \multicolumn{1}{c}{\scriptsize 4C} & \multicolumn{1}{c}{\scriptsize 4F} & \multicolumn{1}{c}{\scriptsize 20} & \multicolumn{1}{c}{\scriptsize 50} & \multicolumn{1}{c}{\scriptsize 52} & \multicolumn{1}{c}{\scriptsize 47} & \multicolumn{1}{c}{\scriptsize 33} & \multicolumn{1}{c}{\scriptsize 32} & \multicolumn{1}{c}{\scriptsize 00}\\
\end{tabular}
\caption{\normalfont Strings in memory; under each cell, the ASCII code in hexadecimal.
Below: \texttt{.asciz "HELLO PRG32"}.}
\label{fig:info:strings}
\end{figure}

The \code{z} of \code{.asciz} means \emph{zero-terminated}: for the eleven
characters of the string of Chapter~\ref{ch:firstcontact} the assembler stores
twelve bytes. The function \code{prg32\_gfx\_text8} receives in \reg{a2} only
the address of the first byte, and reads until it finds the zero. If the zero
were missing, the function would go on reading whatever follows in memory.

\subsection{Unicode and UTF-8: every writing system}

One hundred and twenty-eight characters are enough for English and for
nothing else. \term{Unicode} assigns a number, called a \term{code point}, to
every character of every script; \term{UTF-8} stores that number in one to
four bytes (Table~\ref{tab:info:utf8}). The encoding is compatible with
ASCII: a plain ASCII text is already valid UTF-8. It is the encoding of the
web and of the source files that the reader is writing.

\begin{table}[htb]
\centering\small
\caption{Unicode code points and their UTF-8 encoding.}\label{tab:info:utf8}
\begin{tabular}{@{}llll@{}}
\toprule
\textbf{Character} & \textbf{Code point} & \textbf{UTF-8 bytes} & \textbf{Length}\\
\midrule
A & U+0041 & \texttt{41} & 1 byte: identical to ASCII\\
\`e & U+00E8 & \texttt{C3 A8} & 2 bytes\\
\texteuro & U+20AC & \texttt{E2 82 AC} & 3 bytes\\
the ideogram for ``sun'', ``day'' & U+65E5 & \texttt{E6 97 A5} & 3 bytes\\
the ``grinning face'' emoji & U+1F600 & \texttt{F0 9F 98 80} & 4 bytes\\
\bottomrule
\end{tabular}
\end{table}

\begin{prgwarn}
In UTF-8 the length of a string in bytes is not its length in characters.
\end{prgwarn}

\subsection{Byte order: little-endian}

A word is four bytes, and memory is addressed byte by byte. In which order
do the four bytes of a word go? The directive \code{.word 0x12345678} can be
stored in two ways (Figure~\ref{fig:info:endian}).

\begin{figure}[htb]
\centering\small
\begin{tabular}{@{}r@{\quad}*{4}{|>{\centering\arraybackslash\ttfamily}p{13mm}}|@{}}
\multicolumn{1}{r}{} & \multicolumn{1}{c}{\scriptsize address $A$} & \multicolumn{1}{c}{\scriptsize $A+1$} & \multicolumn{1}{c}{\scriptsize $A+2$} & \multicolumn{1}{c}{\scriptsize $A+3$}\\\cline{2-5}
\textbf{little-endian} (\riscv, x86, ARM) & \cellcolor{prggreen!25}78 & \cellcolor{prgviolet!20}56 & \cellcolor{prgamber!25}34 & \cellcolor{prgaccent!20}12\\\cline{2-5}
\multicolumn{5}{c}{}\\[-6pt]\cline{2-5}
\textbf{big-endian} (network protocols) & \cellcolor{prgaccent!20}12 & \cellcolor{prgamber!25}34 & \cellcolor{prgviolet!20}56 & \cellcolor{prggreen!25}78\\\cline{2-5}
\end{tabular}
\caption{The two byte orders of the word \hex{12345678}.}
\label{fig:info:endian}
\end{figure}

\begin{definition}{Little-endian and big-endian}
In \term{little-endian} order the least significant byte of a word is stored
at the lowest address; in \term{big-endian} order, the most significant byte.
\end{definition}

\riscv is little-endian: this is what \ins{lw} and \ins{sw} do. The byte
order is noticed in two circumstances: when memory is examined in a debugger,
and when data are exchanged with another machine. The names come from
\emph{Gulliver's Travels}, where two peoples go to war over the end from which
an egg should be broken.

% ----------------------------------------------------------------------------
\section{Data structures}

From single values one passes to organised collections. Five ways of
organising data recur in all of computing.

\begin{description}[style=nextline,leftmargin=1.5em]
\item[Array.] Elements of the same type, one after another, reached by index:
  32, 11, 34, 22.
\item[List.] Elements linked one to the next: cheap to insert and to remove,
  slow to reach by position.
\item[Stack.] Last in, first out: \emph{push} adds on top, \emph{pop} removes
  from the top. The frame that saved \reg{ra} in Chapter~\ref{ch:firstcontact}
  was on a stack.
\item[Dictionary.] Pairs of key and value: \code{NA} gives \code{NAPOLI},
  \code{PA} gives \code{PALERMO}, \code{MI} gives \code{MILANO}.
\item[Tree.] A hierarchy: an animal walks or swims, has lungs or gills, two
  legs or four.
\end{description}

In the machine all of them are the same thing: bytes in memory, plus a rule
for finding the next element.

\subsection{An array in memory, an element in a register}

For the array the rule is arithmetic (Listing~\ref{lst:info:array}). An array
is a base address followed by elements of equal size; element $i$ of an array
of words is at address $\mathit{base}+4i$, because a word is four bytes.

\asmfile[caption={Reading \code{scores[2]}.},label=lst:info:array]{code/l04/s19.S}

The instruction \ins{slli} shifts the bits of a register to the left; a shift
by two places multiplies by four, as a shift by one place in decimal
multiplies by ten. (The property is proved in Chapter~\ref{ch:integers}.)
The expression \code{scores[i]} of a high-level language is exactly these four
instructions.

\begin{prgwarn}
What does \ins{lw} read if $i=4$? Whatever follows the array in memory. The
machine does not check bounds: the programmer does.
\end{prgwarn}

% ----------------------------------------------------------------------------
\section{Data and algorithms}

The chapter closes by widening the view, from the bits of one program to the
data of the world and to the procedures that work on them.

\subsection{Data, big data and data science}

\begin{definition}{Big data}
\term{Big data} are data that exceed the limits of traditional database tools.
\end{definition}

They are structured and unstructured: the economy, geolocation, time series,
social networks. Their size is measured in petabytes and zettabytes, that is,
thousands of millions of terabytes. They come from four kinds of source:
\emph{sensor data} (the environment, GPS positions, the Internet of Things:
every device measures something); \emph{transactions} (purchases at the
supermarket, credit cards, online orders); \emph{public data} (administrative
records, territory, consumption: the open data of institutions); and
\emph{social media} (preferences, politics, sentiment: what people write and
share). \term{Data science} is the analysis of large quantities of
information, and the technologies that extract knowledge and value from data
are parallel computing and GPUs.

\begin{definition}{Data mining}
\term{Data mining} is the set of techniques and methodologies for extracting
information or knowledge from large quantities of data.
\end{definition}

Among its techniques are decision trees and classification, rule learning and
Bayesian methods, neural networks and pattern recognition.

\subsection{GPGPU: general-purpose computing on graphics processors}

A graphics processor (GPU) has thousands of simple cores built to do the same
operation on many data at once. The sum of $N$ items, which takes $N$ steps
one after another, needs only $\log_2 N$ parallel steps
(Figure~\ref{fig:info:gpgpu}): the items are added in pairs, then the partial
sums in pairs, and so on. The same hardware that draws video games now trains
neural networks. Under every core there are still registers, instructions and
bits; Chapter~\ref{ch:beyond} returns to the subject.

\begin{figure}[htb]
\centering
\includegraphics[width=0.6\textwidth]{gpgpu_sum.png}
\caption{A parallel sum on the cores of a GPU: $\log_2 N$ steps.}
\label{fig:info:gpgpu}
\end{figure}

\subsection{Algorithms and computability}

A problem is \term{computable} if it can be solved by an algorithm. The
\emph{Turing machine} is an ideal computer with a program, expressed
mathematically through sets, relations and functions. The \emph{von Neumann
machine} is the computing architecture that is able to execute a program which
realises an algorithm for the Turing machine. A problem that no algorithm on
a Turing machine can solve cannot be solved by any computer: this statement,
in its informal version, is the Church--Turing thesis. Not every problem has
an algorithm; the classic example of a problem that is not computable is the
\emph{halting problem}, that of deciding whether an arbitrary program will
eventually stop.

\begin{definition}{Algorithm}
An algorithm is a procedure with five properties. It is \textbf{atomic}: its
steps are elementary and cannot be broken down any further. It is
\textbf{unambiguous}: each step is interpreted in one direct way by the
executor, human or machine. It is \textbf{finite}: it has a finite number of
steps and a finite quantity of input data. It is \textbf{terminating}: its
execution ends after a finite time. It is \textbf{effective}: its execution
leads to a single, well-defined result.
\end{definition}

A \riscv instruction is the model of an atomic and unambiguous step: its
meaning is fixed by the specification, bit by bit.

\subsection{Deterministic and non-deterministic algorithms}

An algorithm is \term{deterministic} if for every instruction there is, with
the same input, exactly one next step: the path of the algorithm exists and is
unique, $y=f(x)$. Converting a unit of measure is deterministic; so is
\code{first\_update}, which always adds one.

An algorithm is \term{non-deterministic} if, with the same input, it can
behave differently at each execution. \emph{Probabilistic} algorithms use a
generator of random numbers; \emph{concurrent} algorithms can depend on timing
(the so-called race conditions). A game that places an enemy by calling
\code{prg32\_random\_number} is non-deterministic.

Listing~\ref{lst:info:prob} shows two probabilistic algorithms in
pseudo-code. The first tests whether a number is prime: the answer
``composite'' is always right, the answer ``probably prime'' is right with
overwhelming probability. Trading a tiny chance of error for a huge gain in
speed is how real cryptography tests primes. (The Fermat test shown is the
simplest; real libraries use the Miller--Rabin test, which closes its
loopholes.) The second is a random walk, which gives a different result at
every run: simulations average many of them.

\txtfile[caption={Two probabilistic algorithms.},label=lst:info:prob]{code/l04/s28.txt}

\subsection{Machine learning}

\term{Machine learning} is a method of data analysis that automates the
construction of an analytical model: the computer finds hidden patterns
without being explicitly programmed where to look. In a classical program the
function $y=f(x)$ is written by a person; in machine learning, $f$ is learned
from examples. Today's generative models are the same idea at an enormous
scale: a large language model predicts the next piece of text from thousands
of millions of learned numbers.

Those numbers are bits in memory, multiplied and added by processors: the
subject of this book. And a model can be wrong with great confidence, so that
the fundamentals are needed to check it. This is the reason for the rule
stated in Chapter~\ref{ch:intro}: the examination verifies that the
fundamentals are in the student's head.

% ----------------------------------------------------------------------------
\section{Check your understanding}

\begin{quiz}{How many different values can one byte represent?}
\item 8
\item 16
\item 255
\item 256
\end{quiz}

\begin{quiz}{The \prg colour \hex{F81F} is:}
\item red
\item magenta
\item cyan
\item yellow
\end{quiz}

\begin{quiz}{How many bytes does \code{.asciz "CIAO"} occupy in memory?}
\item 3
\item 4
\item 5
\item 8
\end{quiz}

\begin{quiz}{On \riscv, \code{.word 0x0000ABCD} is stored at address $A$. Which byte is at address $A$?}
\item \hex{00}
\item \hex{AB}
\item \hex{CD}
\item \hex{DC}
\end{quiz}

\section{Exercises}

\begin{exercise}[on \prg]
\begin{enumerate}[nosep]
\item Change the bar of Chapter~\ref{ch:firstcontact} to magenta, then to a
  dark grey of your choice: compute the RGB565 value by hand first.
\item Print your name with \code{prg32\_gfx\_text8}. How many bytes does the
  string occupy?
\item Declare an array of four colours with \code{.word} and draw four bars,
  one per colour, reading each colour with \ins{lw}.
\item In \prgqt open the memory view on your string and on your array: find
  the terminating zero and check the byte order.
\end{enumerate}
\expected{Four bars in four colours; in the memory view the string ends with
\texttt{00} and each \code{.word} shows its least significant byte first. A
dark grey, for comparison: red 8 of 31, green 16 of 63, blue 8 of 31 is
\texttt{01000\,010000\,01000}, that is, \hex{4208}.}
\end{exercise}

\section{Glossary}

\begin{glossarytable}
bit / byte / word & bit / byte / parola\\
positional notation & notazione posizionale\\
hexadecimal & esadecimale\\
sampling & campionamento\\
quantisation & quantizzazione\\
byte order (endianness) & ordinamento dei byte\\
string & stringa\\
array & vettore (array)\\
linked list & lista concatenata\\
stack & pila\\
tree & albero\\
machine learning & apprendimento automatico\\
\end{glossarytable}

\section{Conclusions and reading}

Digital information is a pattern of bits plus a convention for reading it.
The same sixteen bits can be a number, a colour or two characters: the program
decides. Hexadecimal is a compact spelling of binary, four bits per digit.
Text is numbers (ASCII, UTF-8); structures are bytes plus a rule for finding
the next element. Data are the currency of this realm, and algorithms turn
data into information fit for decisions.

\paragraph{Reading.} \citet{patterson2020}, sections 2.4 and~2.9; Charles
Petzold, \emph{Code}, the chapters on bits, bytes and text; for machine
learning, Sebastian Raschka, \emph{Machine Learning with Python}.
